
\subsubsection{Main Results}

Table 1 presents our cost-performance comparison across 6 UQ method variants on TruthfulQA selective prediction.

\textbf{Table 1: Cost-Performance Comparison of UQ Methods}

\begin{table*}[htbp]
\centering
\resizebox{\dimexpr 2\columnwidth+\columnsep\relax}{!}{%
\begin{tabular}{llllll}
\toprule
Method & Cost & AUROC & Spearman $\rho$ & Gate ($\geq 0.70)$ & Pareto-Optimal \\
\midrule
Temperature Scaling & 1.$0\times$ & 0.682 & 0.26 & No & Yes \\
Conformal Prediction & 1.$0\times$ & 0.695 & 0.23 & No & Yes \\
MC Dropout k=1 & 1.$0\times$ & 0.678 & 0.21 & No & No (dominated) \\
MC Dropout k=3 & 3.$0\times$ & 0.704 & 0.31 & Yes & Yes \\
MC Dropout k=5 & 5.$0\times$ & 0.712 & 0.36 & Yes & Yes \\
MC Dropout k=10 & 10.$0\times$ & 0.718 & 0.38 & Yes & Yes \\
\bottomrule
\end{tabular}
}
\end{table*}

\textbf{Observations:}

Five methods are Pareto-optimal (temperature scaling, conformal prediction, MC dropout k=3/5/10). MC dropout k=1 is dominated (same $1\times$ cost as temperature scaling but lower AUROC 0.678 versus 0.682). This confirms that cost-performance trade-offs exist and no universal method dominates across budget constraints.

MC dropout $k\geq 3$ required for AUROC $\geq$ 0.70 threshold. Zero-cost methods (temperature scaling 0.682, conformal prediction 0.695) and MC dropout k=1 (0.678) fall below the threshold. MC dropout k=3 (0.704) is the minimum configuration achieving the threshold, establishing an epistemic uncertainty requirement at 8B scale.

MC dropout k=10 achieves highest AUROC (0.718) but only marginally outperforms k=5 (0.712, difference 0.006). This reveals diminishing returns: doubling cost from k=5 to k=10 yields negligible AUROC gain.

Temperature scaling not competitive with MC dropout k=5 (|0.682 - 0.712| = 0.030 > threshold). Zero-cost methods trade 3-4\% AUROC for computational savings.

All methods exceed Spearman $\rho$ > 0.2 threshold, confirming uncertainty scores correlate with incorrectness. MC dropout k=10 shows strongest correlation ($\rho$=0.38), temperature scaling weakest ($\rho$=0.26).

\subsubsection{Pareto Frontier Analysis}

The Pareto frontier reveals three distinct efficiency zones:

\textbf{$1\times$ Cost Zone:} Temperature scaling (0.682) versus conformal prediction (0.695). Both Pareto-optimal despite falling below threshold. Conformal achieves +0.013 AUROC. Suitable for budget-constrained applications tolerating lower uncertainty quality.

\textbf{$3\times$ Cost Zone:} MC dropout k=3 (0.704) is an efficiency point. Achieves threshold. Optimal for applications requiring $\geq 0.70$ AUROC with moderate budget constraints.

\textbf{$5\times$ Cost Zone:} MC dropout k=5 (0.712) offers +0.008 AUROC over k=3. Suitable for applications justifying $5\times$ cost for marginal quality improvement.

\textbf{$10\times$ Cost Zone:} MC dropout k=10 (0.718) shows diminishing returns (+0.006 AUROC versus k=5). Exceeds practical budget constraints.

\subsubsection{MC Dropout k-Dependency}

AUROC increases monotonically with k: k=1 (0.678) $\rightarrow$ k=3 (0.704) $\rightarrow$ k=5 (0.712) $\rightarrow$ k=10 (0.718). Marginal gains decrease: k=$1\rightarrow k=3$ (+0.026), k=$3\rightarrow k=5$ (+0.008), k=$5\rightarrow k=10$ (+0.006). This diminishing returns pattern suggests Bayesian approximation converges at $k\approx 5$ for 8B models.

\subsubsection{Baseline Accuracy}

Llama-3.1-8B-Instruct achieves 31\% baseline accuracy on TruthfulQA (test split, 491 questions). MC dropout $k\geq 3$ achieves AUROC $\geq$ 0.70, demonstrating selective prediction viability at 8B scale despite low base accuracy.
